\chapter{Architecture Generality and Comparison with Other Methods}
\label{ch:generality}

The results so far concern one predictor, REL\_GLOBAL, and its ADS-24
adaptation. This chapter asks two further questions. The first is
whether encoders commonly used for log sequences are unstable under
unseen conditions in the same way, and whether CD, frozen as it is,
reduces that instability for them as well
(Section~\ref{sec:gen-arch}). The second is whether simpler
single-model methods achieve the same reduction
(Section~\ref{sec:gen-methods}). Both studies cover the in-house lab
and ADS-24, and the status of the evidence differs between the two
datasets; each section opens by stating it.

\section{Generality across encoders}
\label{sec:gen-arch}

The in-house results of this section are sealed criteria read on 156
new runs (1,600 rows) collected for this study in the PAG1 round. They
repeat the laboratory, benign workers, and configuration design of the
earlier rounds. The ADS-24 results reuse the 35 environment-evaluation
runs scored in Chapter~\ref{ch:ads24} and are therefore development
comparisons (Section~\ref{sec:eval-rounds}). The three encoders are the
LSTM encoder (DeepLog-style)~\cite{hochreiter1997long, du2017deeplog},
the Transformer encoder (LogBERT-style)~\cite{vaswani2017attention,
guo2021logbert}, and the plain GRU~\cite{cho2014gru}. As
Section~\ref{sec:method-transfer} describes, they are encoder
adaptations to our supervised task, not reproductions of DeepLog or
LogBERT: LogitNorm, input-gradient regularization, and the relational
residual are all removed, and width and learning rate were chosen by
validation inside FIT alone. In stage~1 we trained 30 baseline seeds of
each encoder on each dataset (180 fits). In stage~2, with the
computation sealed in stage~1, we trained a CD student for every seed
of the six encoder--dataset combinations, using five teachers chosen
from the same encoder's baselines by unlabeled agreement, and a
self-distillation control for the in-house plain GRU (210 fits).
% 근거: RECORD:279--281; ORIG:ARCH_GENERALITY_STAGE2_20261122/
% ARCH_GENERALITY_STAGE2_REPORT.md:3 (ADS-24 ON35 = 이미 열람한 평가 실행의 개발 비교,
% PAG1 새 156실행 = 같은 lab·worker·구성 설계의 확인), :20 (1단계 seed SD ≥ .020인
% 여섯 조합만 실행), :161 (학생 fit 210/210); ARCH_GENERALITY_STAGE1_REPORT.md:13
% (PROPHET 총합 160계획, 156적격, 1,600행). 1단계 180 fit(+24 pilot)은 표 tab:rounds.
% 원 파일 이름 대응: DEEPLOG_LSTM = LSTM encoder (DeepLog-style),
% LOGBERT_ENCODER = Transformer encoder (LogBERT-style), PLAIN_GRU = plain GRU,
% REL_BASE = baseline, KD_ENV = CD student, KD_SELF = self-distillation control.

All three encoders varied across seeds on the conditions absent from
FIT (Table~\ref{tab:gen-arch}). On \texttt{A+C FLAT} the seed standard
deviation of NLL was 0.216 for the LSTM encoder, 0.833 for the
Transformer encoder, and 0.579 for the plain GRU, against 0.080 for
REL\_GLOBAL on the same runs, and the accuracy of plain-GRU seeds ranged
from 51.0\% to 87.7\%. We had expected that encoders without LogitNorm,
which keeps logits at a fixed norm, would spread less. The opposite
happened. The Transformer figure is driven partly by one seed whose
\texttt{A+C FLAT} loss was 4.690. Without it, the remaining 29 seeds
still had a standard deviation of 0.466 on the A+C score, against 0.040
for their CD students. On ADS-24 the spread was smaller in absolute
terms, but the environment evaluation again separated the seeds far
more than the preservation runs did.
% 근거: RECORD:287, 415 (예상 철회); ARCH_GENERALITY_STAGE2_REPORT.md "전체 성능과
% 정확도" (PLAIN_GRU FLAT 정확도 0.510264–0.876833; LOGBERT FLAT 최악 seed
% 4500015/4.689571), :131 (4500015 제외 민감도: A+C 1:1 평균 기준 0.918233/0.533915,
% SD 0.466200/0.039555; 주 판정은 30개로만 수행). ADS 보존 OFF SD 0.040/0.017/0.012
% 대 환경 ON 0.352/0.126/0.180 (같은 표).

\begin{table}[htbp]
\centering
\caption{Baseline and CD student for the three general encoders
(30 seeds each). In-house values are on \texttt{A+C FLAT} of the new
PAG1 runs (31 runs, 341 rows); ADS-24 values are on the 35
environment-evaluation runs (5,639 rows), a development comparison.
Each cell gives baseline $\rightarrow$ CD student. The last two rows
are the predictors of Chapters~\ref{ch:results} and~\ref{ch:ads24} on
the same runs, for reference (in-house: REL\_GLOBAL $\rightarrow$
KD\_ALL, frozen v3 checkpoints, Table~\ref{tab:results-pag1}); they are
not part of the six-combination family.}
\label{tab:gen-arch}
\small
\begin{tabular}{@{}llccc@{}}
\toprule
Data & Encoder & Mean NLL & Seed SD & Accuracy range (\%) \\
\midrule
In-house & LSTM (DeepLog-style)
  & 0.651 $\rightarrow$ 0.486 & 0.216 $\rightarrow$ 0.050
  & 65.7--88.3 $\rightarrow$ 80.9--86.2 \\
In-house & Transformer (LogBERT-style)
  & 1.050 $\rightarrow$ 0.548 & 0.833 $\rightarrow$ 0.040
  & 28.4--83.3 $\rightarrow$ 73.0--79.5 \\
In-house & plain GRU
  & 1.139 $\rightarrow$ 0.258 & 0.579 $\rightarrow$ 0.015
  & 51.0--87.7 $\rightarrow$ 84.2--89.1 \\
ADS-24 & LSTM (DeepLog-style)
  & 2.249 $\rightarrow$ 2.079 & 0.352 $\rightarrow$ 0.038
  & 22.3--32.7 $\rightarrow$ 29.3--31.1 \\
ADS-24 & Transformer (LogBERT-style)
  & 2.134 $\rightarrow$ 1.796 & 0.126 $\rightarrow$ 0.027
  & 26.3--32.6 $\rightarrow$ 32.9--34.3 \\
ADS-24 & plain GRU
  & 2.221 $\rightarrow$ 2.130 & 0.180 $\rightarrow$ 0.026
  & 24.6--33.2 $\rightarrow$ 29.2--31.6 \\
\midrule
In-house & REL\_GLOBAL (reference)
  & 0.134 $\rightarrow$ 0.078 & 0.080 $\rightarrow$ 0.024
  & 89.1--100.0 $\rightarrow$ 93.8--99.7 \\
ADS-24 & REL\_GLOBAL (reference)
  & 2.276 $\rightarrow$ 2.060 & 0.163 $\rightarrow$ 0.025 & --- \\
\bottomrule
\end{tabular}
\end{table}
% 근거: RECORD:283--292; ARCH_GENERALITY_STAGE2_REPORT.md "전체 성능과 정확도"
% (EVAL_ON, A_PLUS_C_FLAT 행), "별도 확인: 동결된 v3 30 seed를 PAG1 새 156실행에
% 재추론" (REL_GLOBAL 0.133967/0.080252 [0.891496–1.000000], KD_ALL 0.078413/0.024293
% [0.938416–0.997067]). ADS-24 REL_GLOBAL 행은 ADS24_STAGE2_REPORT.md (RECORD:218--219).
% 정확도는 원 파일 값을 소수 첫째 자리 %로 줄였다. 기록 표의 정수 % 범위와 같은 값이다.
% ADS-24 REL_GLOBAL의 seed 정확도 범위는 기록에 없어 비웠다.

CD lowered both the spread and the mean in all six combinations, and
all six passed the three sealed criteria of stability, mean improvement,
and preservation under Bonferroni intervals adjusted for six contrasts
(99.1667\%). The closest stability result was the ADS-24 Transformer
encoder, with a seed standard deviation difference of $-0.099$
$[-0.127, -0.075]$; the closest mean result was the ADS-24 plain GRU, at
$-0.092$ $[-0.113, -0.071]$. Every interval for preservation also lay
below its limit, although one only narrowly, as discussed below. The
students' spread was also small relative to the baselines' spread: in
every combination, even the worst CD seed scored better than the median
baseline seed.
% 근거: RECORD:279, 294; ARCH_GENERALITY_STAGE2_REPORT.md "주 판정" (99.1667% 구간:
% ADS LOGBERT SD −0.099170 [−0.127209, −0.075001]; ADS PLAIN_GRU 평균 −0.091707
% [−0.112726, −0.071334]; 여섯 조합 모두 Y·Y·Y). PROPHET 대비는 A+C FLAT/NESTED 1:1
% 평균이다(같은 파일 :76).
% TODO(사람): "가장 나쁜 CD seed가 기준 모델의 중간 seed보다 낫다"(RECORD:279)는
% 기록의 문장입니다. seed별 CSV(STAGE2_SEED_RESULTS.csv)가 atlas에 없어 중간값을
% 원 파일에서 대조하지 못했습니다. 확인해 주세요.

Not every seed improved. Depending on the combination, between 20 and
30 of the 30 seeds had a lower loss after distillation than before.
Seeds that had been better than the consensus of their teachers moved
toward it and gave up part of their advantage, so the best baseline
seed was sometimes better than every student. On ADS-24, for instance, two LSTM baseline seeds
scored below the majority-class baseline and none of the LSTM students
did. What CD delivers is a predictable level of performance, close to
that of the teacher average, rather than the best level any seed could
reach. A practitioner who could identify the good seeds in advance
would lose by distilling them; Chapters~\ref{ch:diagnosis}
and~\ref{ch:ads24} showed that the seen conditions do not identify
them.
% 근거: RECORD:297, 424; ARCH_GENERALITY_STAGE2_REPORT.md "교사 수준, 단순 기준선,
% 민감도" (ADS DEEPLOG_LSTM 다수 기준보다 나은 seed 기준 2/30, KD 0/30; 학생−교사
% NLL 0.009~0.050). 20~30개의 조합별 값은 기록에만 있고 원 파일에서 찾지 못했다.

On ADS-24, lower spread did not mean a useful absolute level. Only the
Transformer students beat the majority-class baseline of 1.830 defined
in Chapter~\ref{ch:ads24}, 27 of 30 seeds with a mean 0.034 below it.
The LSTM and plain-GRU students, like REL\_GLOBAL in
Chapter~\ref{ch:ads24}, remained above the baseline with every seed, by
0.250 and 0.300 on average. In the in-house lab every student of every
encoder was below the corresponding baseline of 1.412. CD narrows the
distribution of seeds; it does not remove the loss caused by the move
to an unseen environment.
% 근거: RECORD:298, 423; ARCH_GENERALITY_STAGE2_REPORT.md:5, "교사 수준, 단순 기준선,
% 민감도" (기준 1.829580; KD 0/30, 27/30, 0/30; 평균 NLL−다수 기준 KD 0.249872,
% −0.034009, 0.300071; PROPHET 기준 1.411751, KD 30/30 세 구조 모두).
% PROPHET 다수 기준(1.411751)의 정의는 원 파일에서 확인하지 못했다.
% TODO(사람): 자체 실험실 다수 클래스 기준선 1.412가 무엇으로 정해졌는지(FIT 다수
% 행동, 어느 평가층) 확인해 주세요. 확인되지 않으면 이 문장을 지우겠습니다.

Two in-house details qualify the result for the plain GRU. Its
preservation interval on \texttt{ORIGINAL\_18} had an upper end of
$+0.0198$, just below the limit of $+0.020$, and its seed-mean accuracy
on that condition fell from 84.5\% to 84.0\%; preservation passed, but
with almost no margin. The self-distillation control, in which each seed
is distilled from its own frozen baseline with the same KL term, scored
1.112 with a seed standard deviation of 0.558 on \texttt{A+C FLAT},
essentially the baseline values of 1.139 and 0.579. This is the second
dataset, after ADS-24, on which distillation without a consensus left
the instability in place. For context, adaptations of three security
systems were run with three seeds each on the same new inputs: ORTHRUS
and BiTA reached an A+C NLL of 0.13 to 0.16, close to REL\_GLOBAL, and
CL-AP$^{2}$ about 0.97~\cite{jiang2025orthrus, nayeri2026bita,
liu2024clap2}. Three seeds are too few to measure their spread, so
these figures say nothing about stability. Two attempts to extend CD,
a confidence-conditioned target and a projection onto a representation
shared by the seeds, failed their preliminary checks and were not
carried forward (Appendix~\ref{app:rejected}).
% 근거: RECORD:296, 299, 301, 302; ARCH_GENERALITY_STAGE2_REPORT.md "주 판정"
% (PROPHET PLAIN_GRU 보존 −0.037245, 99.1667% [−0.096992, 0.019818]),
% "전체 성능과 정확도" (ORIGINAL_18 정확도 0.845410 → 0.839917; KD_SELF FLAT
% 1.112231/0.557598), :133 (KD_SELF는 주 판정 가족 밖의 보조 대조).
% +0.0198은 개요의 넷째 자리 예외다.
% ORIG:ARCH_GENERALITY_STAGE1_20261122/SECURITY_PAPER_ADDENDUM/REPORT.md
% (C-fold FIT603, seed 11/29/47, 동결 checkpoint를 PAG1 적격 실행에 적용;
% A+C FLAT/NESTED: CL-AP² 0.967384/0.975969, ORTHRUS 0.145795/0.133481,
% BiTA 0.159013/0.132978; ADS-24 적용판 없음; 원 논문 파이프라인 재현 아님).
% TODO(사람): 부록 C(c_rejected.tex)의 label을 app:rejected로 가정했습니다.

\section{Comparison with other single-model methods}
\label{sec:gen-methods}

The second study compares CD with five other ways of training one model,
all on the plain GRU. The in-house runs are again newly collected under
a sealed contract in the same lab design; the ADS-24 runs are the
inspected environment-evaluation runs, so that half is a development
comparison. The contract fixed the contrasts and intervals but no
pass-or-fail rule (Table~\ref{tab:rounds}, note~c). Weight averaging
(SWA) averages the weights of the baseline's own training trajectory
at updates 1,010, 1,122, 1,234, and 1,346~\cite{izmailov2018averaging}.
Pseudo-labeling adds a cross-entropy term toward the hard predictions
of the same seed's frozen baseline on the unlabeled environment
inputs~\cite{tmp_lee}. Entropy minimization adds the entropy of the
student's own predictions on those inputs~\cite{tmp_grandvalet}.
FIT-input distillation follows the same five-teacher average as CD but
on FIT inputs instead of environment inputs. Geometric-pool distillation
follows the normalized geometric mean of the five teachers on
environment inputs~\cite{tmp_ilichev}. Every method kept the FIT loss,
the 1,346 updates, the coefficient of 0.25 on the added term, a
temperature of 1, and 32 added rows per update, and the 30 seeds started
from the same initial weights. The study trained 300 new models; the
baseline, CD, and self-distillation groups were reused from
Section~\ref{sec:gen-arch}.
% 근거: RECORD:338, 452; ORIG:METHOD_BASELINES_20261123/METHOD_BASELINES_REPORT.md:3
% (ADS-24 ON35 = 이미 열람한 실행의 사후 방법 비교, PROPHET = 같은 실험실 설계의
% 새 실행); METHOD_BASELINES_CONTRACT_SEALED.md:7 (다섯 기준선, 300 fit), :15
% (1,346 update × 32 FIT occurrence, 같은 초기 tensor), :17 (계수 0.25, 온도 1, 32행),
% :19--27 (군별 계산; PSEUDO 신뢰도 문턱 없음; FIT_KD 교사 FIT 출력은 in-sample;
% BASE·KD_ENV·KD_SELF 재사용), :35--37 (판정 규칙 없이 네 통계와 99% 구간).
% 원 파일 이름 대응: SWA = weight averaging, PSEUDO = pseudo-labeling,
% ENTROPY = entropy minimization, FIT_KD = FIT-input distillation,
% GEO_KD = geometric-pool distillation, KD_ENV = CD student.
% TODO(사람): Lee(pseudo-label), Grandvalet·Bengio(entropy minimization)는 2장과 같은
% 임시 key(tmp_lee, tmp_grandvalet)로 인용했습니다. 기하평균 교사의 출처로 개요가
% 지목한 Ilichev 등(RANLP 2021)은 임시 key tmp_ilichev로 인용했습니다.

Of the six training methods, only the two that follow a consensus of
several models on environment inputs reduced the spread
(Table~\ref{tab:gen-methods}). On \texttt{A+C FLAT} the CD student had
a seed standard deviation of 0.013 and the geometric-pool student 0.020,
while the baseline and the other five groups stayed between 0.32 and
1.22. On
ADS-24 the pattern was the same, with 0.026 and 0.030 against 0.17 to
0.49. Table~\ref{tab:gen-methods-contrast} gives the paired contrasts.
In the in-house lab CD had a lower mean NLL than weight averaging,
pseudo-labeling, entropy minimization, and FIT-input distillation, with
every interval above zero. Against geometric-pool distillation the
direction was not resolved: the mean difference was $-0.011$
$[-0.067, 0.027]$ and the difference in seed standard deviation $0.003$
$[-0.003, 0.018]$. On ADS-24, CD had a lower mean than all five methods,
including the geometric pool. Even so, its mean of 2.130 remained above
the majority-class baseline of 1.830, and of all 210 models in the
comparison on ADS-24 only one seed of FIT-input distillation fell below
that baseline.
% 근거: RECORD:340--349, 356; METHOD_BASELINES_REPORT.md:7--8, "주 비교: 단일 기준선
% − KD_ENV" (ADS24, PROPHET 표), :37 (다수 기준보다 나은 seed: BASE 0, KD_ENV 0,
% SWA 0, PSEUDO 0, ENTROPY 0, FIT_KD 1, GEO_KD 0; 7군 × 30 = 210), :39, :53.
% 이 값들은 개요 질문 A1에 따라 원 파일에서 옮겼다(RECORD:358, 502는 "미확인"으로
% 적었음).

\begin{table}[htbp]
\centering
\caption{Point estimates for the method comparison on the plain GRU
(30 seeds per method). In-house columns are \texttt{A+C FLAT} of the
new runs of this comparison (21 runs, 231 rows); ADS-24 columns are the
35 environment-evaluation runs (5,639 rows), a development comparison.
The in-house runs differ from those of Table~\ref{tab:gen-arch}, so the
baseline and CD values of the two tables are not compared with each
other. Self-distillation was not part of the ADS-24 comparison.}
\label{tab:gen-methods}
\small
\begin{tabular}{@{}lrrcrr@{}}
\toprule
 & \multicolumn{3}{c}{In-house \texttt{A+C FLAT}}
 & \multicolumn{2}{c}{ADS-24} \\
\cmidrule(lr){2-4}\cmidrule(l){5-6}
Method & NLL & Seed SD & Accuracy range (\%) & NLL & Seed SD \\
\midrule
Baseline                      & 1.233 & 0.693 & 45.9--93.1 & 2.221 & 0.180 \\
CD student                    & 0.222 & 0.013 & 87.9--93.1 & 2.130 & 0.026 \\
Geometric-pool distillation   & 0.216 & 0.020 & 87.4--93.1 & 2.238 & 0.030 \\
Weight averaging (SWA)        & 1.212 & 0.681 & 46.3--89.2 & 2.201 & 0.174 \\
FIT-input distillation        & 1.290 & 0.815 & 43.7--90.0 & 2.226 & 0.193 \\
Self-distillation             & 1.171 & 0.691 & 44.6--93.9 & ---   & ---   \\
Pseudo-labeling               & 2.012 & 1.215 & 44.2--90.5 & 3.355 & 0.438 \\
Entropy minimization          & 2.702 & 0.324 & 80.5 (all seeds) & 6.873 & 0.490 \\
\bottomrule
\end{tabular}
\end{table}
% 근거: RECORD:340--349; METHOD_BASELINES_REPORT.md "전 평가층의 전체 성과"
% (ADS24 EVAL_ON, PROPHET A_PLUS_C_FLAT 21/231). 정확도 범위는 원 파일 값을 소수
% 첫째 자리 %로 줄였다(예: BASE 0.458874–0.930736). 기록 표의 정수 % 범위와 같다.

\begin{table}[htbp]
\centering
\caption{Paired contrasts of the method comparison, each method minus
the CD student, with two-sided Bonferroni intervals at 99\% for the
five methods of each dataset. Positive values favor CD. The in-house
score here is the 1:1 average of \texttt{A+C FLAT} and
\texttt{A+C NESTED}, not the \texttt{A+C FLAT} value of
Table~\ref{tab:gen-methods}; point estimates use all eligible runs, and
intervals resample the configurations with both repeats complete. The
last column states whether the upper end of the method's preservation
interval against the baseline lies below the limit ($+0.020$ in-house,
$+0.043$ on ADS-24).}
\label{tab:gen-methods-contrast}
\small
\begin{tabular}{@{}llrlrlc@{}}
\toprule
Data & Method & \multicolumn{2}{c}{Seed SD difference}
 & \multicolumn{2}{c}{Mean NLL difference} & Pres. \\
\midrule
In-house & Weight averaging       & 0.664 & $[0.593, 0.745]$ & 0.985 & $[0.875, 1.085]$ & yes \\
In-house & Pseudo-labeling        & 1.194 & $[1.103, 1.329]$ & 1.789 & $[1.602, 1.973]$ & no \\
In-house & Entropy minimization   & 0.291 & $[0.207, 0.444]$ & 2.374 & $[1.703, 3.408]$ & no \\
In-house & FIT-input distillation & 0.802 & $[0.730, 0.883]$ & 1.065 & $[0.921, 1.206]$ & yes \\
In-house & Geometric-pool distillation & 0.003 & $[-0.003, 0.018]$ & $-0.011$ & $[-0.067, 0.027]$ & yes \\
\midrule
ADS-24 & Weight averaging       & 0.148 & $[0.116, 0.180]$ & 0.071 & $[0.054, 0.089]$ & yes \\
ADS-24 & Pseudo-labeling        & 0.412 & $[0.345, 0.482]$ & 1.225 & $[1.004, 1.437]$ & no \\
ADS-24 & Entropy minimization   & 0.463 & $[0.394, 0.541]$ & 4.743 & $[4.088, 5.347]$ & no \\
ADS-24 & FIT-input distillation & 0.166 & $[0.131, 0.200]$ & 0.097 & $[0.075, 0.120]$ & yes \\
ADS-24 & Geometric-pool distillation & 0.004 & $[0.002, 0.006]$ & 0.108 & $[0.082, 0.133]$ & yes \\
\bottomrule
\end{tabular}
\end{table}
% 근거: METHOD_BASELINES_REPORT.md "주 비교: 단일 기준선 − KD_ENV" (:29--35, :43--49),
% :51 (PROPHET 점추정은 전체 적격 실행, 구간은 두 반복이 완성된 구성);
% METHOD_BASELINES_CONTRACT_SEALED.md:33 (PROPHET 주 점수 = FLAT·NESTED 1:1 평균),
% :35 (보존은 각 기준선 − BASE, ADS 3%, PROPHET +.020), :37 (100,000회 PCG64,
% 고정 30 seed). 보존 "no"는 원 파일의 "미확인/초과"(상단이 한도를 넘음)다.
% 정확도 SD 차이와 최저 seed 정확도 차이는 분량 때문에 표에서 뺐다.
% 개요 질문 A1: 이 판정 구간은 원 파일에 있으므로 본문에 썼다.

The controls separate three parts of the method. A consensus of several
models is needed: following one's own frozen copy, softly in
self-distillation or by hard labels in pseudo-labeling, either changed
nothing or made the loss worse. The environment inputs are needed:
following the same five-teacher average on FIT inputs left the spread
at the baseline level, because on FIT the teachers largely agree and
the student has little to learn from them. And an average along a
single training trajectory is not enough: weight averaging stayed at
the baseline's spread on both datasets. Between the arithmetic and the
geometric pool of the same teachers we cannot choose. They were
indistinguishable in the in-house lab, and the arithmetic pool was
better on ADS-24 only in a development comparison; CD keeps the
arithmetic mean of Section~\ref{sec:method-cd}. Entropy minimization is
a warning about reading low disagreement as success. It made all 30
seeds agree, so that every seed had the same accuracy of 80.5\% on
\texttt{A+C FLAT}, while their loss rose to 2.702.
Chapter~\ref{ch:decomposition} returns to this case.
% 근거: RECORD:351--356, 404 (기하평균 교사 교체: 교사 NLL 56회 중 24회만 우세,
% 학생 우열 미확인, 산술평균 유지); METHOD_BASELINES_REPORT.md "V와 합의 예측"
% (ENTROPY A_PLUS_C_FLAT V 0.000007, 합의 NLL 2.701564; 이 수치는 10장에서 쓴다).
% "FIT에서 교사들이 대체로 같다": FIT view top-1 완전 일치 90.3%(74,074/82,008),
% 개요 B 대조와 11.4절. 이 값은 자체 모델(REL_GLOBAL) 교사에 대한 것이며 plain GRU
% 교사에 대해서는 확인하지 않았다.
% TODO(사람): 위의 "on FIT the teachers largely agree" 구절은 REL_GLOBAL 교사의
% FIT 일치도(90.3%)에 근거한 해석입니다. plain GRU 교사에 대한 값이 없으므로, 이
% 구절을 빼거나 근거를 알려 주세요.

This comparison has several limits. The in-house collection was planned
as 160 attempts, and a sealed safety stop ended it after 119, leaving
110 eligible runs and 1,119 rows; the 41 runs not attempted were neither
retried nor replaced. The \texttt{A+C} intervals therefore rest on six
complete \texttt{FLAT} configurations and nine complete \texttt{NESTED}
configurations. The comparison used one encoder, and no method was tuned
on its own: coefficient, exposure, and schedule were fixed for all, and
weight averaging did not use the cyclical or constant learning-rate
schedule of the original SWA. In-house, the weight-averaging contrast may
also include rerun variation, because the plain GRU on CPU is not
deterministic: the cross-entropy run from which the weights were
averaged ended with weights identical, bit for bit, to the stored
baseline for 30 of 30 seeds on ADS-24 and for none in-house.
Finally, the plain-GRU CD student is not a substitute for REL\_GLOBAL.
On the same new runs the A+C score of REL\_GLOBAL was 0.086, against
0.226 for the plain-GRU student.
% 근거: RECORD:358, 452, 504; METHOD_BASELINES_REPORT.md:9 (REL_GLOBAL 참고 0.085753,
% KD_ENV 0.226057, A+C 1:1), :10, :15--16 (계획 160, 시도 119, THREE_SOURCE_FAILURES,
% 미시도 41, 완전 구성 FLAT 6/16, NESTED 9/16), :18 (SWA bitwise 일치 ADS 30/30,
% PROPHET 0/30, CPU GRU 비결정성).

% TODO(사람): 온도 보정과 확률 앙상블(개요 질문 A1). 원 파일에 값이 있지만 기록은
% 마무리 확인 묶음 목록(RECORD:498)에 올려 두었으므로 기본값대로 본문에 쓰지 않았다.
% 본문에 써도 되면 아래 값으로 한 문단을 추가하겠다.
% METHOD_BASELINES_REPORT.md "앙상블·온도 보정 참고" (못 본 조건 BASE / 상위 5 앙상블 /
% 30 앙상블 / 온도 보정 seed 평균; PROPHET은 A+C 1:1):
%   ADS24 LSTM 2.248586 / 2.053577 / 1.841886 / 2.229396
%   ADS24 Transformer 2.134342 / 1.783663 / 1.753357 / 2.064998
%   ADS24 plain GRU 2.221359 / 2.079862 / 1.962718 / 2.163016
%   ADS24 REL_GLOBAL 2.276041 / 2.037885 / 1.945412 / 2.184814
%   PROPHET LSTM 0.695279 / 0.427706 / 0.320941 / 0.763100
%   PROPHET Transformer 0.993020 / 0.414734 / 0.421735 / 0.919611
%   PROPHET plain GRU 1.236530 / 0.213731 / 0.445200 / 1.219041
%   PROPHET REL_GLOBAL 0.085753 / 0.022455 / 0.048526 / 0.102500
%   (PROPHET REL_GLOBAL 온도는 FIT 120실행 in-sample 참고치.)
% 온도 보정 − CD, 여섯 대비 99.1667% (SINGLE_MODEL_SCOPE_CLARIFICATION.md: 봉인 뒤,
% 점수 열람 전에 범위 추가): ADS24 평균 0.033364 [0.016628, 0.049431], SD 0.137918
% [0.107162, 0.167891]; PROPHET 평균 0.992984 [0.885849, 1.101081], SD 0.666835
% [0.593450, 0.751775].
